\sectionheading{5.38 Loop — ধারণা ও প্রয়োজনীয়তা}

Program-এ অনেক সময় একই কাজ বারবার করতে হয়।
একই Statement বারবার আলাদাভাবে না লিখে নির্দিষ্ট নিয়মে
পুনরাবৃত্তি করার জন্য \textbf{Loop} ব্যবহার করা হয়।

%================================================
% Definition
%================================================

\booksubheading{Loop কী?}

\begin{definitionbox}

কোনো নির্দিষ্ট Condition বা Sequence-এর ভিত্তিতে একই Statement
বা Code Block বারবার Execute করার প্রক্রিয়াকে
\textbf{Loop} বা \textbf{Iteration} বলে।

\end{definitionbox}

\begin{center}

\textbf{একই কাজ}
\(\longrightarrow\)
\textbf{বারবার}
\(\longrightarrow\)
\textbf{Loop}

\end{center}

%================================================
% Without loop vs loop
%================================================

\noindent
\begin{minipage}[t]{0.49\linewidth}

\begin{sideexamplebox}{Loop ছাড়া}

১ থেকে ৫ পর্যন্ত সংখ্যা Print করতে—

\begin{sidecodebox}
\begin{english}
\ttfamily
print(1)\\
print(2)\\
print(3)\\
print(4)\\
print(5)
\end{english}
\end{sidecodebox}

একই ধরনের Statement বারবার লিখতে হয়েছে।

\end{sideexamplebox}

\end{minipage}
\hfill
\begin{minipage}[t]{0.49\linewidth}

\begin{sideexamplebox}{Loop ব্যবহার করে}

\begin{sidecodebox}
\begin{english}
\ttfamily
for i in range(1, 6):\\
\hspace*{4mm}print(i)
\end{english}
\end{sidecodebox}

\textbf{Output:}

\begin{center}
\textenglish{\ttfamily 1 \quad 2 \quad 3 \quad 4 \quad 5}
\end{center}

কম Code-এ একই কাজ করা যায়।

\end{sideexamplebox}

\end{minipage}

%================================================
% Why Loop
%================================================

\booksubheading{Loop কেন প্রয়োজন?}

Loop ব্যবহারের প্রধান প্রয়োজন—

\begin{compacttwocolpoints}

\item একই কাজ বারবার করার জন্য।

\item একই Code বহুবার লেখা এড়াতে।

\item Program ছোট, পরিষ্কার ও সহজ রাখতে।

\item Code-এর Repetition কমাতে।

\item List, String বা অন্যান্য Collection-এর Item Process করতে।

\item Counting, Summation, Searching ইত্যাদি কাজে।

\item নির্দিষ্ট Condition পূরণ হওয়া পর্যন্ত কাজ চালিয়ে যেতে।

\end{compacttwocolpoints}

\begin{rememberbox}

\begin{center}

\textbf{Loop}
\(\longrightarrow\)
\textbf{Repetition কমায়}
\(\longrightarrow\)
\textbf{Code ছোট ও সহজ হয়}

\end{center}

\end{rememberbox}

%================================================
% Iteration + Loop Body
%================================================

\booksubheading{Iteration ও Loop Body}

\noindent
\begin{minipage}[t]{0.49\linewidth}

\begin{sideexamplebox}{Iteration}

Loop-এর প্রতিবার Code Block Execute হওয়াকে
\textbf{Iteration} বলে।

যদি কোনো Loop ৫ বার Execute হয়, তাহলে সেখানে
\textbf{৫টি Iteration} হয়েছে।

\medskip

\textbf{MCQ তথ্য:}

\begin{center}
\textbf{Loop Body-এর একবার Execution}
\(\longrightarrow\)
\textbf{এক Iteration}
\end{center}

\end{sideexamplebox}

\end{minipage}
\hfill
\begin{minipage}[t]{0.49\linewidth}

\begin{sideexamplebox}{Loop Body}

Loop-এর যে Statement বা Code Block বারবার Execute হয় তাকে
\textbf{Loop Body} বলে।

\begin{sidecodebox}
\begin{english}
\ttfamily
for i in range(3):\\
\hspace*{4mm}print("ICT")
\end{english}
\end{sidecodebox}

এখানে
\textenglish{\ttfamily print("ICT")}
হলো Loop Body এবং এটি ৩ বার Execute হবে।

\end{sideexamplebox}

\end{minipage}

%================================================
% General Flow
%================================================

\booksubheading{Loop-এর সাধারণ কার্যধারা}

সাধারণভাবে Loop-এর Logic—

\begin{center}

\textbf{Start}
\(\longrightarrow\)
\textbf{Condition / Next Item}
\(\longrightarrow\)
\textbf{Loop Body}

\medskip

\textbf{Loop Body}
\(\longrightarrow\)
\textbf{Update / Next Item}
\(\longrightarrow\)
\textbf{আবার পরীক্ষা}
\(\longrightarrow\)
\textbf{Stop}

\end{center}

Loop-এর ধরন অনুযায়ী বাস্তব গঠন কিছুটা ভিন্ন হতে পারে।

%================================================
% Important Parts
%================================================

\booksubheading{Loop-এর গুরুত্বপূর্ণ অংশ}

\begin{tabularx}
{\linewidth}
{
|>{\bfseries\RaggedRight\arraybackslash}p{0.27\linewidth}
|Y|
}

\hline
\rowcolor{softBlue}

\multicolumn{1}{|c|}{\textbf{বিষয়}}
&
\multicolumn{1}{c|}{\textbf{কাজ}}
\\
\hline

Initialization
&
শুরুতে প্রয়োজনীয় Variable-এর প্রাথমিক Value নির্ধারণ।
\\
\hline

Condition / Range
&
Loop চলবে কি না বা কতদূর চলবে তা নির্ধারণ।
\\
\hline

Loop Body
&
বারবার Execute হওয়া Statement বা Code Block।
\\
\hline

Update
&
প্রয়োজন অনুযায়ী Counter/Variable-এর Value পরিবর্তন।
\\
\hline

Iteration
&
Loop Body-এর প্রতিবার Execution।
\\
\hline

\end{tabularx}

%================================================
% while anatomy
%================================================

\begin{examplebox}{\textenglish{\ttfamily while} Loop-এ অংশগুলো চেনো}

\begin{codebox}
\begin{english}
\ttfamily
i = 1\\[0.5mm]
while i <= 5:\\
\hspace*{4mm}print(i)\\
\hspace*{4mm}i += 1
\end{english}
\end{codebox}

\begin{center}

\textenglish{\ttfamily i = 1}
\(\longrightarrow\)
\textbf{Initialization}

\qquad

\textenglish{\ttfamily i <= 5}
\(\longrightarrow\)
\textbf{Condition}

\medskip

\textenglish{\ttfamily print(i)}
\(\longrightarrow\)
\textbf{Loop Body}

\qquad

\textenglish{\ttfamily i += 1}
\(\longrightarrow\)
\textbf{Update}

\end{center}

\end{examplebox}

%================================================
% Counter
%================================================

\booksubheading{Counter বা Loop Variable}

Loop কতবার চলছে বা কোন Item নিয়ে কাজ করছে তা নির্দেশ বা
নিয়ন্ত্রণ করতে ব্যবহৃত Variable-কে অনেক ক্ষেত্রে
\textbf{Counter} বা \textbf{Loop Variable} বলা হয়।

\noindent
\begin{minipage}[t]{0.49\linewidth}

\begin{sideexamplebox}{\textenglish{\ttfamily for}-এ Loop Variable}

\begin{sidecodebox}
\begin{english}
\ttfamily
for i in range(5):\\
\hspace*{4mm}print(i)
\end{english}
\end{sidecodebox}

এখানে
\textenglish{\ttfamily i}
হলো Loop Variable।

\end{sideexamplebox}

\end{minipage}
\hfill
\begin{minipage}[t]{0.49\linewidth}

\begin{sideexamplebox}{\textenglish{\ttfamily while}-এ Counter}

\begin{sidecodebox}
\begin{english}
\ttfamily
i = 1\\
while i <= 5:\\
\hspace*{4mm}print(i)\\
\hspace*{4mm}i += 1
\end{english}
\end{sidecodebox}

এখানে
\textenglish{\ttfamily i}
Counter হিসেবে কাজ করছে।

\end{sideexamplebox}

\end{minipage}

%================================================
% Types of Loop
%================================================

\booksubheading{Python-এর প্রধান Loop}

Python-এ প্রধানত দুই ধরনের Loop ব্যবহৃত হয়—

\begin{tabularx}
{\linewidth}
{
|>{\bfseries\RaggedRight\arraybackslash}p{0.28\linewidth}
|Y|
}

\hline
\rowcolor{softBlue}

\multicolumn{1}{|c|}{\textbf{Loop}}
&
\multicolumn{1}{c|}{\textbf{সাধারণ ব্যবহার}}
\\
\hline

\textenglish{\ttfamily for} Loop
&
Sequence/Iterable-এর Item নিয়ে কাজ বা নির্দিষ্ট সংখ্যক Iteration।
\\
\hline

\textenglish{\ttfamily while} Loop
&
কোনো Condition True থাকা পর্যন্ত Repetition।
\\
\hline

\end{tabularx}

\begin{rememberbox}

\begin{center}

\textenglish{\ttfamily for}
\(\longrightarrow\)
\textbf{Sequence / নির্দিষ্ট Iteration}

\qquad

\textenglish{\ttfamily while}
\(\longrightarrow\)
\textbf{Condition-নির্ভর Repetition}

\end{center}

\end{rememberbox}

%================================================
% for + while basic examples
%================================================

\noindent
\begin{minipage}[t]{0.49\linewidth}

\begin{sideexamplebox}{\textenglish{\ttfamily for} Loop-এর ধারণা}

\textenglish{\ttfamily for}
সাধারণত কোনো Sequence বা Iterable-এর প্রতিটি Item নিয়ে
একবার করে কাজ করে।

\begin{sidecodebox}
\begin{english}
\ttfamily
for x in [10, 20, 30]:\\
\hspace*{4mm}print(x)
\end{english}
\end{sidecodebox}

\textbf{Output:}

\begin{center}
\textenglish{\ttfamily 10 \quad 20 \quad 30}
\end{center}

\end{sideexamplebox}

\end{minipage}
\hfill
\begin{minipage}[t]{0.49\linewidth}

\begin{sideexamplebox}{\textenglish{\ttfamily while} Loop-এর ধারণা}

Condition যতক্ষণ True থাকে ততক্ষণ Code Block পুনরাবৃত্তি করে।

\begin{sidecodebox}
\begin{english}
\ttfamily
i = 1\\
while i <= 3:\\
\hspace*{4mm}print(i)\\
\hspace*{4mm}i += 1
\end{english}
\end{sidecodebox}

\textbf{Output:}

\begin{center}
\textenglish{\ttfamily 1 \quad 2 \quad 3}
\end{center}

\end{sideexamplebox}

\end{minipage}

%================================================
% For vs While
%================================================

\booksubheading{\textenglish{\ttfamily for} ও \textenglish{\ttfamily while}-এর প্রাথমিক পার্থক্য}

\begin{tabularx}
{\linewidth}
{
|>{\bfseries\RaggedRight\arraybackslash}p{0.23\linewidth}
|Y
|Y|
}

\hline
\rowcolor{softBlue}

\multicolumn{1}{|c|}{\textbf{বিষয়}}
&
\multicolumn{1}{c|}{\textenglish{\ttfamily for}}
&
\multicolumn{1}{c|}{\textenglish{\ttfamily while}}
\\
\hline

মূল ভিত্তি
&
Sequence/Iterable বা Iteration
&
Condition
\\
\hline

সাধারণ ব্যবহার
&
Item বা Iteration সম্পর্কে ধারণা থাকলে সুবিধাজনক
&
Condition-এর উপর Repetition নির্ভর করলে সুবিধাজনক
\\
\hline

Variable Control
&
Loop Variable সাধারণত Item গ্রহণ করে
&
প্রয়োজনে Variable Initialize/Update করতে হয়
\\
\hline

Infinite Loop-এর ঝুঁকি
&
তুলনামূলক কম
&
Condition/Update ভুল হলে বেশি
\\
\hline

\end{tabularx}

%================================================
% Finite + Infinite
%================================================

\booksubheading{Finite Loop ও Infinite Loop}

\noindent
\begin{minipage}[t]{0.49\linewidth}

\begin{sideexamplebox}{Finite Loop}

যে Loop নির্দিষ্ট Iteration শেষে বন্ধ হয় তাকে
\textbf{Finite Loop} বলা যায়।

\begin{sidecodebox}
\begin{english}
\ttfamily
for i in range(5):\\
\hspace*{4mm}print(i)
\end{english}
\end{sidecodebox}

এটি ৫ Iteration শেষে শেষ হবে।

\end{sideexamplebox}

\end{minipage}
\hfill
\begin{minipage}[t]{0.49\linewidth}

\begin{sideexamplebox}{Infinite Loop}

যে Loop-এর Ending Condition কার্যকরভাবে কখনো False হয় না,
সে Loop চলতেই থাকে।

\begin{sidecodebox}
\begin{english}
\ttfamily
while True:\\
\hspace*{4mm}print("Hello")
\end{english}
\end{sidecodebox}

এটি নিজে থেকে বন্ধ হবে না।

\end{sideexamplebox}

\end{minipage}

%================================================
% Common infinite loop mistake
%================================================

\begin{rememberbox}

\textbf{\color{deepOrange}ভুল করো না}

\begin{codebox}
\begin{english}
\ttfamily
i = 1\\
while i <= 5:\\
\hspace*{4mm}print(i)
\end{english}
\end{codebox}

এখানে
\textenglish{\ttfamily i}
Update করা হয়নি। তাই
\textenglish{\ttfamily i <= 5}
Condition সবসময় True থাকবে এবং Loop চলতেই থাকবে।

\begin{center}

\textenglish{\ttfamily while}
Loop-এ প্রয়োজনীয়
\textbf{Update ভুলে গেলে}
\(\longrightarrow\)
\textbf{Infinite Loop হতে পারে}

\end{center}

\end{rememberbox}

%================================================
% Loop Condition
%================================================

\booksubheading{Loop Condition}

Loop কখন চলবে এবং কখন শেষ হবে তা নির্ধারণে Condition গুরুত্বপূর্ণ।

\begin{codebox}
\begin{english}
\ttfamily
while i <= 5:
\end{english}
\end{codebox}

\begin{center}

\textbf{Condition True}
\(\longrightarrow\)
\textbf{Loop Body Execute}

\qquad

\textbf{Condition False}
\(\longrightarrow\)
\textbf{Loop শেষ}

\end{center}

\begin{rememberbox}

\begin{center}

\textenglish{\ttfamily True}
\(\longrightarrow\)
\textbf{আবার Loop}

\qquad

\textenglish{\ttfamily False}
\(\longrightarrow\)
\textbf{Loop থেকে বের হও}

\end{center}

\end{rememberbox}

%================================================
% Benefits
%================================================

\booksubheading{Loop-এর সুবিধা}

\begin{compacttwocolpoints}

\item একই Statement বারবার লিখতে হয় না।

\item Code Length কমে এবং Repetition কমে।

\item Program সহজে পরিবর্তন ও Maintain করা যায়।

\item Repetitive Task দ্রুত করা যায়।

\item Collection-এর একাধিক Item Process করা সহজ হয়।

\item Counting, Searching, Summation ইত্যাদি সহজে করা যায়।

\item Program বেশি Organized ও Readable হয়।

\end{compacttwocolpoints}

%================================================
% Same output comparison
%================================================

\booksubheading{Loop ছাড়া বনাম Loop ব্যবহার}

ধরা যাক,
\textenglish{\ttfamily "Python"}
পাঁচবার Print করতে হবে।

\noindent
\begin{minipage}[t]{0.49\linewidth}

\begin{sideexamplebox}{Loop ছাড়া}

\begin{sidecodebox}
\begin{english}
\ttfamily
print("Python")\\
print("Python")\\
print("Python")\\
print("Python")\\
print("Python")
\end{english}
\end{sidecodebox}

একই Statement ৫ বার লেখা হয়েছে।

\end{sideexamplebox}

\end{minipage}
\hfill
\begin{minipage}[t]{0.49\linewidth}

\begin{sideexamplebox}{Loop ব্যবহার করে}

\begin{sidecodebox}
\begin{english}
\ttfamily
for i in range(5):\\
\hspace*{4mm}print("Python")
\end{english}
\end{sidecodebox}

Output একই, কিন্তু Code অনেক সংক্ষিপ্ত।

\end{sideexamplebox}

\end{minipage}

%================================================
% Nested Loop
%================================================

\booksubheading{Nested Loop-এর প্রাথমিক ধারণা}

একটি Loop-এর ভিতরে আরেকটি Loop থাকলে তাকে
\textbf{Nested Loop} বলে।

\begin{examplebox}{উদাহরণ}

\begin{codebox}
\begin{english}
\ttfamily
for i in range(2):\\
\hspace*{4mm}for j in range(3):\\
\hspace*{8mm}print(i, j)
\end{english}
\end{codebox}

Outer Loop-এর প্রতিটি Iteration-এর জন্য Inner Loop সম্পূর্ণভাবে
Execute হয়।

\end{examplebox}

\begin{rememberbox}

Nested Loop পরবর্তী
\textbf{Loop Pattern}
অংশে বিশেষভাবে কাজে লাগবে।

\end{rememberbox}

%================================================
% Uses of Loop
%================================================

\booksubheading{Loop কোথায় ব্যবহার হয়?}

\begin{tabularx}
{\linewidth}
{
|Y
|Y|
}

\hline
\rowcolor{softBlue}

\multicolumn{1}{|c|}{\textbf{ব্যবহার}}
&
\multicolumn{1}{c|}{\textbf{ব্যবহার}}
\\
\hline

১ থেকে
\textenglish{\ttfamily n}
পর্যন্ত সংখ্যা Print
&
জোড়/বিজোড় সংখ্যা বের করা
\\
\hline

Sum বা Average নির্ণয়
&
Multiplication Table তৈরি
\\
\hline

List/String-এর Item Process
&
Searching ও Counting
\\
\hline

Pattern তৈরি
&
Condition পূরণ হওয়া পর্যন্ত Input গ্রহণ
\\
\hline

\end{tabularx}

%================================================
% break continue preview
%================================================

\booksubheading{\textenglish{\ttfamily break} ও \textenglish{\ttfamily continue}-এর প্রাথমিক ধারণা}

Loop-এর স্বাভাবিক Execution পরিবর্তনের জন্য Python-এ
\textenglish{\ttfamily break}
ও
\textenglish{\ttfamily continue}
ব্যবহার করা যায়।

\begin{tabularx}
{\linewidth}
{
|>{\bfseries\RaggedRight\arraybackslash}p{0.28\linewidth}
|Y|
}

\hline
\rowcolor{softBlue}

\multicolumn{1}{|c|}{\textbf{Statement}}
&
\multicolumn{1}{c|}{\textbf{প্রাথমিক কাজ}}
\\
\hline

\textenglish{\ttfamily break}
&
Loop সঙ্গে সঙ্গে বন্ধ করে।
\\
\hline

\textenglish{\ttfamily continue}
&
বর্তমান Iteration-এর বাকি অংশ Skip করে পরবর্তী Iteration-এ যায়।
\\
\hline

\end{tabularx}

\begin{rememberbox}

এগুলো বিস্তারিতভাবে
\textbf{5.41 Loop Control Statement}
অংশে আলোচনা করা হবে।

\begin{center}

\textenglish{\ttfamily break}
\(\longrightarrow\)
\textbf{Loop বন্ধ}

\qquad

\textenglish{\ttfamily continue}
\(\longrightarrow\)
\textbf{Current Iteration Skip}

\end{center}

\end{rememberbox}

%================================================
% Common Mistakes
%================================================

\booksubheading{সাধারণ ভুল}

\begin{rememberbox}

\begin{itemize}

\item Loop Body-তে সঠিক Indentation না দেওয়া।

\item
\textenglish{\ttfamily while}
-এর Condition ভুল লেখা।

\item
\textenglish{\ttfamily while}
Loop-এ প্রয়োজনীয় Counter Update না করা।

\item প্রয়োজন ছাড়া Infinite Loop তৈরি করা।

\item Loop কতবার Iterate করবে তা ভুল হিসাব করা।

\item
\textenglish{\ttfamily for}
ও
\textenglish{\ttfamily while}
-এর ব্যবহার গুলিয়ে ফেলা।

\item Loop Variable-এর Value কীভাবে পরিবর্তিত হচ্ছে তা ভুল বোঝা।

\item Nested Loop-এ Outer ও Inner Loop গুলিয়ে ফেলা।

\end{itemize}

\end{rememberbox}

%================================================
% Final Reminder
%================================================

\begin{rememberbox}

\textbf{\color{deepOrange}মনে রাখবে}

\begin{tabularx}
{\linewidth}
{
>{\bfseries\RaggedRight\arraybackslash}p{0.24\linewidth}
Y
>{\bfseries\RaggedRight\arraybackslash}p{0.24\linewidth}
Y
}

Loop
&
একই কাজ বারবার
&
Iteration
&
Loop Body-এর একবার Execution
\\[0.7mm]

Loop Body
&
পুনরাবৃত্ত Code
&
Counter
&
Iteration নিয়ন্ত্রণ/গণনার Variable
\\[0.7mm]

\textenglish{\ttfamily for}
&
Sequence/Iteration ভিত্তিক
&
\textenglish{\ttfamily while}
&
Condition ভিত্তিক
\\[0.7mm]

Finite Loop
&
নির্দিষ্ট সময়ে শেষ
&
Infinite Loop
&
অনির্দিষ্টকাল চলতে পারে
\\[0.7mm]

Nested Loop
&
Loop-এর ভিতরে Loop
&
Update
&
Counter/Variable পরিবর্তন
\\[0.7mm]

\textenglish{\ttfamily break}
&
Loop বন্ধ
&
\textenglish{\ttfamily continue}
&
Current Iteration Skip
\\

\end{tabularx}

\medskip

\begin{center}

\textbf{Start}
\(\longrightarrow\)
\textbf{Condition / Next Item}
\(\longrightarrow\)
\textbf{Loop Body}

\medskip

\textbf{Loop Body}
\(\longrightarrow\)
\textbf{Update / Next Item}
\(\longrightarrow\)
\textbf{Repeat}
\(\longrightarrow\)
\textbf{Stop}

\medskip

\textenglish{\ttfamily for}
\(\longrightarrow\)
\textbf{Sequence / Iteration}

\qquad

\textenglish{\ttfamily while}
\(\longrightarrow\)
\textbf{Condition}

\end{center}

\end{rememberbox}

\small

%================================================
% 5.39 Python for Loop
%================================================

